\documentclass[10pt,a4paper]{amsart}
\usepackage[margin=19mm]{geometry}
\usepackage{amsmath,amssymb,mathtools}
\usepackage[scr=rsfs,cal=euler]{mathalfa}
\usepackage{xfrac}
\usepackage[unicode,hidelinks,bookmarks=false]{hyperref}
\newcommand{\ie}{i.e.\ }
\newcommand{\cf}{cf.\ }
\newcommand{\resp}{respectively\ }
\newcommand{\Subsection}{\subsection}
\providecommand{\nobreakdash}{\nobreak\hbox{-}}
\setlength{\emergencystretch}{2em}
\multlinegap0pt
\usepackage{amssymb}
\newtheorem{lemma}{Lemma}
\title{The Diophantine equation $p^{x}+6^{y}=z^{2n}$: a complete solution for half of the primes}
\author{Pagdame Tiebekabe}
\date{Source: arXiv:2609.09214v1 (CC BY 4.0)}
\begin{document}
\maketitle

\noindent\emph{Source and licence.} The Diophantine equation $p^{x}+6^{y}=z^{2n}$: a complete solution for half of the primes. Version arXiv:2609.09214v1 (CC BY 4.0), distributed by arXiv under the Creative Commons Attribution 4.0 International licence, http://creativecommons.org/licenses/by/4.0/, as recorded in the arXiv licence metadata for this exact version and shown on the abstract page https://arxiv.org/abs/2609.09214v1. Full licence notice: \texttt{licenses/arxiv-2609.09214-CC-BY-4.0.txt}.\par\smallskip

\noindent\textit{Editorial scope.} Counted unit: the article's lemma lem:consecutive (source0.tex lines 261--266). The article's complete original proof of it is delivered with it (source0.tex lines 268--284, one \texttt{proof} environment of 17 lines). No mathematics is written, repaired or re-derived: the statement and the proof are the article's own lines, in the article's own notation, and no retained line is substituted or re-flowed.  No further source-local context is needed: every cross-reference of every retained line resolves inside the two delivered spans.  Nothing was omitted from any retained span, so the package carries no omission record.  No retained line contains a citation, so the wrapper carries no bibliography and no bibliography entry.  No retained line contains a figure environment, an image-inclusion command or drawing code, so no image is delivered and the package declares no asset dependency.  Editorial formatting only, in addition: an AMS \texttt{amsart} preamble that restates the article's own declarations and macros used by the retained lines, namely the environments \texttt{lemma} on separate counters, together with the standard packages those lines need.  The retained environments print in this document as Lemma 1 in order of appearance, whereas the article prints the same items with the numbering of its own full document.  The complete native source of the article as distributed in the arXiv e-print for this exact version is bundled unchanged as \texttt{sources/209790/source0.tex}.
\begin{lemma}\label{lem:consecutive}
There is no pair $(T,m)$ of integers with $T\geq 2$ and $m\geq 1$ such that
\[
(T-1)(T+1)=2^{\,m+1}\,3^{m}.
\]
\end{lemma}

\begin{proof}
Write $T-1=2u$ and $T+1=2v$, so that $v-u=1$, $\gcd(u,v)=1$ and
$uv=2^{\,m-1}3^{m}$. Since $u$ and $v$ are coprime, each is composed of a
single prime, and the condition $v-u=1$ leaves the three possibilities
\[
\{u,v\}=\bigl\{1,\,2^{\,m-1}3^{m}\bigr\},\quad
\{u,v\}=\bigl\{2^{\,m-1},\,3^{m}\bigr\},\quad
\{u,v\}=\bigl\{3^{m},\,2^{\,m-1}\bigr\}.
\]
The first gives $2^{\,m-1}3^{m}=2$, impossible since $3\mid 2^{\,m-1}3^{m}$.
The second gives $3^{m}-2^{\,m-1}=1$; but $3^{m}-2^{\,m-1}=2$ for $m=1$ and
\[
3^{m}-2^{\,m-1}>3^{m}-3^{m-1}=2\cdot 3^{m-1}\geq 6 \qquad(m\geq 2),
\]
so it is never equal to $1$. The third gives $2^{\,m-1}-3^{m}=1$, which is
impossible since $3^{m}>2^{\,m-1}$.
\end{proof}

\end{document}
